\documentclass[a4paper]{article}
\usepackage[pdftex]{graphicx}
\usepackage[utf8]{inputenc}
\usepackage{enumerate}
\usepackage{siunitx}
\sisetup{locale=DE} 
\usepackage{href-ul}
\hypersetup{
	colorlinks=true,
	linkcolor=blue,
	urlcolor=blue}
\usepackage{icomma}
\usepackage{amssymb}
\usepackage{geometry}
\geometry{a4paper, top=15mm, left=15mm, right=15mm, bottom=15mm,
	headsep=10mm, footskip=12mm}
%Source: img_0101

% Start the document
\begin{document}
%Titel
{\bf Schulaufgabe aus der Physik }\\
\begin{enumerate}[1.]

%Aufgabe 1
{\bf \item In einem gemeinsamen Stromkreis des $\SI{230}{\volt}$--Netzes, der mit $\SI{10}{\ampere}$ abgesichert ist, sind folgende Geräte (parallel) angeschlossen:}\\
 $G_1$ mit $ \SI{75}{\watt};\quad  G_2$ mit $\SI{1,5}{\ampere}; \quad G_3 $ mit $\SI{45}{\ohm}$.

\begin{enumerate}[a)]
\item Berechnen Sie folgende elektrische Größen der angeschlossenen Geräte:
$$ I_1=\dots \quad P_2=\dots \quad I_3=\dots$$
\item Berechnen Sie die maximale Leistung eines weiteren Gerätes $G_4$, das noch zusätzlich in diesen Stromkreis parallel geschaltet werden kann.
\end{enumerate}

%Aufgabe 2
{\bf \item Zwei elektrische Widerstände $R_1$ und $R_2$ werden von einem Strom $I = \SI{0,18}{\ampere}$ durchflossen.} An $R_1$ fällt eine Spannung von $U_1=\SI{9,0}{\volt}$ ab, die Größe $R_2$ ist 
$\SI{150}{\ohm}$.\\
 Berechnen Sie:
$$R_1=\dots \quad U_2=\dots \quad U_{ges}=\dots $$

%Aufgabe 3
{\bf \item Ein elektrisches Gerät mit ($\SI{110}{\volt} / \SI{0,4}{\ampere} $) soll an die Spannung $\SI{230}{\volt}$ ange\-schlossen werden.}  Berechnen Sie die Größe des nötigen elektrischen Vorwiderstandes $R_{vor}$.

%Aufgabe 4
{\bf \item Zeigen Sie mit einer kurzen Rechnung, wie sich der elektrische Widerstand $R$ eines Metalldrahtes verändert, wenn $\dots$}
\begin{enumerate}[a)]
\item $\dots$die Drahtlänge verdoppelt und gleichzeitig die Querschnittsfläche verdreifacht wird.
\item $\dots$die Drahtlänge vervierfacht und der Radius verdoppelt wird
\end{enumerate}

%Aufgabe 5
{\bf \item Eine 17,5 m lange Kupfer-Doppelleitung 
$\bigl( \rho_{Cu}= \SI{0,017}{\ohm\cdot\milli\meter^2 \per \meter}\bigr)$
besitzt einen elektrischen Widerstand von $R=\SI{0,20}{\ohm}$.}\\
Wie groß ist die Querschnittsfläche und der Radius des verwendeten Drahtes?

{\bf Kreuzen Sie in den Aufgaben 6 und 7 alle wahren Aussagen an!}

%Aufgabe 6
\item Bei einer reinen Parallelschaltung von mehreren elektrischen Widerständen gilt
 {\bf immer}:\\
Mit jedem weiteren zugeschalteten Widerstand wird$\dots$\\

\begin{tabular}{p{75 pt}p{75 pt}p{75 pt}p{75 pt}}
$\square R_{ges}$ größer& $\square R_{ges}$ kleiner&
 $\square I_{ges}$ größer&$\square I_{ges}$ kleiner
\end{tabular}

%Aufgabe 7
\item Sind zwei Widerstände $R_1$ und $R_2$ parallel geschaltet, wobei $R_1>R_2$ gilt, so gilt {\bf immer} $\dots$ :
\begin{enumerate}
\item $\dots$ für $R_{ges}$:

\begin{tabular}{p{75 pt}p{75 pt}p{75 pt}p{75 pt}}
  $\square R_{ges}>R_1 $&$ \square R_{ges}<R_1 $&$
 \square R_{ges}>R_2 $&$ \square R_{ges}<R_2 $
\end{tabular}

\item $\dots$  für die Stromstärken $I_1$ und $I_2$:

\begin{tabular}{p{100 pt}p{100 pt}p{100 pt}}
  $\square I_1>I_2 $ & $ \square I_1=I_2 $ & $ \square I_1 <I_2 $
\end{tabular}

\item $\dots$  für die Spannungen  $U_1$ und $U_2$:

\begin{tabular}{p{100 pt}p{100 pt}p{100 pt}}
  $\square U_1>U_2 $ & $ \square U_1=U_2 $ & $ \square U_1 <U_2 $
\end{tabular}
\end{enumerate}

%Aufgabe 8
{\bf \item Welche von den Lämpchen }
$$L_1 ( \SI{4,0}{\volt} / \SI{0,1}{\ampere}); \quad L_2 (\SI{12}{\volt} / \SI{1,2}{\watt})
 ~\textnormal{und} \quad L_3 (\SI{4,0}{\volt} / \SI{20}{\ohm})$$
 können physikalisch sinnvoll in einer Reihenschaltung zusammengeschaltet werden?
\begin{enumerate}[a)]
\item Ermitteln Sie aus den Daten die nötigen physikalischen Größen und entscheiden Sie!
\item Berechnen Sie für die möglichen Lämpchen aus 8.a) die Betriebsstromstärke $I_{ges}$, den Gesamtwiderstand $R_{ges}$, und die Gesamtspannung $U_{ges}$. 
\end{enumerate}

%Aufgabe 9
{\bf \item Die Widerstände $R_1=\SI{100}{\ohm}$ und $R_2=\SI{300}{\ohm}$ sind in Parallelschaltung an die Spannung von $U_{ges}=\SI{150}{\volt}$ angeschlossen.}\\
Berechnen Sie $R_{ges}$.
 
%Aufgabe 10
{\bf \item Berechnen Sie im folgenden Schaltbild$\dots$}

\begin{minipage}{0.65\textwidth}
\includegraphics[width=6 cm]{widerstand101.pdf}
\end{minipage}
\hfill
\begin{minipage}{0.3\textwidth}
\renewcommand{\arraystretch}{3}
\begin{tabular}{cc}
$U_2=\dots$ & $ U_{ges}=\dots$ \\
 $I_{ges}=\dots$ & $ I_1=\dots$ \\
$I_3=\dots$ & $ R_{ges}=\dots$ 
\end{tabular}
\end{minipage}

\end{enumerate}

\fbox{
	\begin{minipage}{0.5\textwidth}
		Zur Lösung bitte \href{https://www.okuyakl.de/phys/p10elL101/ll101.pdf}{hier klicken} oder den QR-Code scannen.\\
	Weitere Arbeitsblätter gibt es unter 
	
	\href{https://www.okuyakl.de}{www.okuyakl.de}
	\end{minipage}
	\hfill
	\begin{minipage}{0.4\textwidth}
		\includegraphics[width=1.5 cm]{../../viecher/zwe03}
		\includegraphics[width=3 cm]{qr101}
		\includegraphics[width=2 cm]{../../viecher/afanticon1}
		
	\end{minipage}}
\end{document}%L Lösung --------------------------------------------------
